% !TEX root = ../00_Main/Main.tex
\documentclass[../00_Main/Main.tex]{subfiles}
\begin{document}

\label{sec:background}

Este cap\'itulo resume los conceptos que se usan en el resto de la tesis. Cubre los procesos de decisi\'on de Markov, los algoritmos de aprendizaje
por refuerzo implementados, la entrop\'ia como medida de confianza, los
ensambles de modelos, la optimizaci\'on bayesiana de hiperpar\'ametros y
las pruebas estad\'isticas con que se comparan los resultados. La
referencia general es el texto de Sutton y Barto~\citep{SuttonBarto2018}.

\subsection{Procesos de decisi\'on de Markov}

Un MDP finito es una tupla
$\mathcal{M} = (\mathcal{S}, \mathcal{A}, P, R, \gamma)$ donde
$\mathcal{S}$ es un espacio de estados finito, $\mathcal{A}$ un
espacio de acciones finito, $P(s' \mid s, a)$ la probabilidad de
transici\'on, $R(s, a, s')$ una funci\'on de recompensa acotada y
$\gamma \in [0, 1)$ el factor de descuento. Bajo una pol\'itica
$\pi(a \mid s)$ el agente genera una trayectoria
$\tau = (s_0, a_0, r_1, s_1, a_1, \ldots)$ cuyo \emph{retorno
descontado} desde el tiempo $t$ es:

\begin{equation}
G_t \;=\; \sum_{k = 0}^{\infty} \gamma^{k}\, r_{t + k + 1}
\label{eq:return}
\end{equation}

Donde $r_{t+k+1}$ es la recompensa recibida $k+1$ pasos despu\'es de $t$.
La funci\'on de valor de estado--acci\'on de la pol\'itica es
$Q^{\pi}(s, a) = \mathbb{E}_{\pi}[G_t \mid s_t = s,\, a_t = a]$ y la
pol\'itica \'optima satisface la ecuaci\'on de Bellman:

\begin{equation}
Q^{*}(s, a) \;=\; \mathbb{E}\!\left[r + \gamma \max_{a'} Q^{*}(s', a')
\,\middle|\, s, a\right]
\label{eq:bellman-opt}
\end{equation}

Donde la esperanza se toma sobre el estado siguiente $s' \sim P(\cdot \mid s, a)$
y la recompensa $r = R(s, a, s')$.

\subsection{Q-Learning tabular}

Cuando $\mathcal{S}$ y $\mathcal{A}$ son enumerables, $Q$ se almacena
como una tabla y se aplica la actualizaci\'on de
Watkins~\citep{Watkins1992}:

\begin{equation}
Q(s, a) \;\leftarrow\; Q(s, a) \;+\; \alpha\,
\bigl[r + \gamma\, \max_{a'} Q(s', a') - Q(s, a)\bigr]
\label{eq:qlearning}
\end{equation}

Donde $\alpha \in (0, 1]$ es la tasa de aprendizaje. Las acciones se
eligen con una pol\'itica $\varepsilon$-\emph{greedy}: con probabilidad
$\varepsilon$ se elige una acci\'on al azar y, en caso contrario,
$\arg\max_a Q(s,a)$.

\subsubsection{Double Q-Learning}

El operador $\max$ de la Ecuaci\'on~\eqref{eq:qlearning} sobreestima el
valor de las acciones. \emph{Double Q-Learning}~\citep{Hasselt2010}
mantiene dos tablas $Q^{A}$ y $Q^{B}$ y usa una para elegir la acci\'on
siguiente y la otra para evaluarla:

\begin{equation}
Q^{A}(s, a) \leftarrow Q^{A}(s, a) + \alpha\bigl[r + \gamma\,
Q^{B}\!\bigl(s', \arg\max_{a'} Q^{A}(s', a')\bigr) - Q^{A}(s, a)\bigr]
\label{eq:double-q}
\end{equation}

La actualizaci\'on de $Q^{B}$ es sim\'etrica, y en cada paso se elige
con probabilidad $1/2$ cu\'al de las dos tablas se actualiza.

\subsubsection{Modelado de recompensa mediante potencial (PBRS)}

El \emph{Potential-Based Reward Shaping} (PBRS)~\citep{Ng1999} suma a
la recompensa el t\'ermino $F(s, s') = \gamma\Phi(s') - \Phi(s)$, donde
$\Phi : \mathcal{S} \to \mathbb{R}$ es una funci\'on de potencial sobre
los estados. Ng \emph{et al.} demostraron que este t\'ermino no cambia
las pol\'iticas \'optimas del MDP original. En Double Q-Learning, el
potencial es la suma de las severidades actuales de las ciudades con
signo negativo, $\Phi(s) = -\sum_i S_i$, y el t\'ermino se multiplica por
un coeficiente $\kappa$: la recompensa efectiva es
$\tilde r = r + \kappa F(s, s')$.

\subsection{Deep Q-Networks (DQN)}

Cuando el espacio de estados es grande, $Q(s, a; \theta)$ se aproxima
con una red neuronal de par\'ametros $\theta$. El \emph{Deep
Q-Network}~\citep{Mnih2015} estabiliza el entrenamiento con dos elementos. El primero es un \emph{buffer} de repetici\'on de experiencia (\emph{experience replay})~\citep{Lin1992}, del que se muestrean lotes de transiciones al azar para romper la correlaci\'on entre muestras consecutivas. El segundo es una \emph{red objetivo} (\emph{target network})
$Q(\cdot;\theta^{-})$, cuyos par\'ametros se copian desde $\theta$ cada
$C$ pasos. La funci\'on de costo es:

\begin{equation}
\mathcal{L}(\theta) \;=\;
\mathbb{E}_{(s,a,r,s') \sim \mathcal{D}}
\!\left[\bigl(r + \gamma \max_{a'} Q(s', a'; \theta^{-})
- Q(s, a; \theta)\bigr)^{2}\right]
\label{eq:dqn-loss}
\end{equation}

Donde $\mathcal{D}$ es el \emph{buffer} de transiciones almacenadas.

\subsubsection{Double DQN}

\emph{Double DQN}~\citep{Hasselt2016} aplica la idea de la
Ecuaci\'on~\eqref{eq:double-q} sin agregar una segunda red: la acci\'on
siguiente se elige con la red en l\'inea $\theta$ y se eval\'ua con la
red objetivo $\theta^{-}$:

\begin{equation}
y \;=\; r + \gamma\, Q\!\bigl(s',\, \arg\max_{a'} Q(s', a'; \theta);\,
\theta^{-}\bigr)
\label{eq:ddqn-target}
\end{equation}

Los tres agentes de tipo DQN de mPES (DQN, DQN recurrente y
Transformer) usan este blanco, enmascaran las acciones no factibles
antes del $\arg\max$ y reemplazan el error cuadr\'atico por la p\'erdida
de Huber~\citep{Huber1964}.

\subsection{Redes recurrentes y LSTM}

Cuando la informaci\'on relevante para decidir no est\'a toda en el
estado actual, es habitual procesar una secuencia de estados con una
red recurrente. La \emph{Long Short-Term Memory}
(LSTM)~\citep{Hochreiter1997} usa compuertas de entrada $i_t$, olvido
$f_t$ y salida $o_t$, con valores en $(0,1)$, que controlan c\'omo el
estado de celda $c_t$ acumula informaci\'on:

\begin{equation}
c_t \;=\; f_t \odot c_{t-1} \;+\; i_t \odot \tilde{c}_t
\qquad h_t \;=\; o_t \odot \tanh(c_t)
\label{eq:lstm}
\end{equation}

Donde $\tilde{c}_t$ es el candidato de celda, $h_t$ el estado oculto y
$\odot$ el producto elemento a elemento. Hausknecht y
Stone~\citep{Hausknecht2015} mostraron que incorporar una capa LSTM en
un DQN (\emph{Deep Recurrent Q-Network}) mejora el desempe\~no en
entornos parcialmente observables.

\subsection{Gradiente de pol\'itica y actor--cr\'itico}

Los m\'etodos de \emph{gradiente de pol\'itica}~\citep{Williams1992}
parametrizan directamente la pol\'itica $\pi_\theta(a \mid s)$ y
ajustan $\theta$ en la direcci\'on que aumenta el retorno esperado
$J(\theta) = \mathbb{E}_{\pi_\theta}[G_0]$:

\begin{equation}
\nabla_{\theta} J(\theta) \;=\;
\mathbb{E}_{\pi_{\theta}}\!\bigl[\nabla_{\theta} \log
\pi_{\theta}(a \mid s)\, A^{\pi_{\theta}}(s, a)\bigr]
\label{eq:pg}
\end{equation}

Donde $A^{\pi}(s, a) = Q^{\pi}(s, a) - V^{\pi}(s)$ es la \emph{ventaja}
y $V^{\pi}(s) = \mathbb{E}_{\pi}[G_t \mid s_t = s]$. El \emph{Advantage
Actor--Critic} (A2C)~\citep{Mnih2016} entrena a la vez un actor, que
produce $\pi_\theta$, y un cr\'itico, que estima $V_{\phi}$. La ventaja
se estima con \emph{Generalized Advantage Estimation}
(GAE)~\citep{Schulman2016}:

\begin{equation}
\hat{A}_t^{\mathrm{GAE}(\gamma,\lambda)} \;=\;
\sum_{k=0}^{\infty}(\gamma\lambda)^{k}\,\delta_{t+k}
\qquad \delta_t = r_{t+1} + \gamma V_\phi(s_{t+1}) - V_\phi(s_t)
\label{eq:gae}
\end{equation}

Donde $\delta_t$ es el error de diferencia temporal de un paso y
$\lambda \in [0,1]$ regula el compromiso entre sesgo y varianza. El actor
se regulariza con un t\'ermino $-\beta_H H(\pi_\theta(\cdot \mid s))$,
con $\beta_H \geq 0$, que evita que la pol\'itica se concentre demasiado
pronto en una acci\'on.

\subsection{Atenci\'on y Transformer}

El \emph{Transformer}~\citep{Vaswani2017} reemplaza la recurrencia por
\emph{atenci\'on de producto escalar escalado}. Dadas las matrices de
consultas $Q$, claves $K$ y valores $V$, obtenidas por proyecciones
lineales de la secuencia de entrada, y siendo $d_k$ la dimensi\'on de
las claves, cada posici\'on combina los valores de las dem\'as con pesos
$\mathrm{softmax}(QK^{\top}/\sqrt{d_k})$. Para que cada posici\'on s\'olo
use informaci\'on pasada se suma una m\'ascara triangular $M$, con
$M_{ij} = -\infty$ para $j > i$ y $M_{ij} = 0$ en otro caso:

\begin{equation}
\mathrm{Atenci\acute{o}n}_{\mathrm{causal}}(Q, K, V) =
\mathrm{softmax}\!\left(\frac{Q K^{\top}}{\sqrt{d_k}} + M\right) V
\label{eq:causal-attention}
\end{equation}

\subsection{Entrop\'ia de Shannon como medida de confianza}

Dada una distribuci\'on de probabilidad discreta $p(a \mid s)$ sobre
las $|\mathcal{A}|$ acciones, su \emph{entrop\'ia de
Shannon}~\citep{Shannon1948} es:

\begin{equation}
H(p) \;=\; -\sum_{a \in \mathcal{A}} p(a \mid s)\, \log p(a \mid s)
\label{eq:shannon}
\end{equation}

Con la convenci\'on $0 \log 0 = 0$. Su valor m\'aximo, $\log
|\mathcal{A}|$, corresponde a la distribuci\'on uniforme, de modo que la
versi\'on normalizada

\begin{equation}
H_{\mathrm{norm}}(p) \;=\; \frac{H(p)}{\log |\mathcal{A}|} \in [0, 1]
\label{eq:shannon-norm}
\end{equation}

vale $0$ cuando toda la probabilidad est\'a en una acci\'on y $1$ cuando
todas las acciones son igualmente probables. La entrop\'ia de la
distribuci\'on de salida se usa para cuantificar la incertidumbre de un
clasificador; por ejemplo, Lakshminarayanan
\emph{et al.}~\citep{Lakshminarayanan2017} la emplean para detectar
entradas ajenas a la distribuci\'on de entrenamiento. En esta tesis se
usa $1 - H_{\mathrm{norm}}$ como \emph{confianza} del agente en cada
decisi\'on. Cuando el modelo produce valores $Q$ y no probabilidades,
primero se convierten en una distribuci\'on (Secci\'on~\ref{sec:ens-methods}), por
lo que la confianza resultante es una medida heur\'istica y no una
probabilidad calibrada.

\subsection{Ensambles de modelos}\label{sec:ensembles-bg}

Un \emph{ensamble} combina las salidas de $K$ modelos ya entrenados
para producir una \'unica decisi\'on. Sea $p_k(a \mid s)$ la
distribuci\'on sobre acciones del miembro $k$. Si el miembro produce
valores $Q_k(s,\cdot)$, la distribuci\'on se obtiene con una
\emph{softmax} de temperatura $\tau$:

\begin{equation}
p_k(a \mid s) \;=\;
\frac{e^{Q_k(s,a)/\tau}}
     {\sum_{a'} e^{Q_k(s,a')/\tau}}
\label{eq:ens-softmax}
\end{equation}

Un $\tau$ peque\~no concentra la distribuci\'on en la acci\'on de mayor
valor y uno grande la acerca a la uniforme. La regla de combinaci\'on
m\'as habitual es el \emph{voto suave} (\emph{soft voting}), un promedio
ponderado de las distribuciones:

\begin{equation}
p_{\mathrm{ens}}(a \mid s) \;=\;
\frac{\sum_{k=1}^{K} w_k \, p_k(a \mid s)}{\sum_{k=1}^{K} w_k}
\qquad w_k \geq 0
\label{eq:soft-voting}
\end{equation}

Donde $w_k$ es el peso del miembro $k$; la acci\'on elegida es
$\arg\max_a p_{\mathrm{ens}}(a \mid s)$. En el \emph{voto duro}
(\emph{hard voting}), en cambio, cada miembro vota s\'olo por su acci\'on
preferida $a_k = \arg\max_a p_k(a \mid s)$ y gana la m\'as votada:

\begin{equation}
a_{\mathrm{ens}} \;=\;
\arg\max_{a} \sum_{k=1}^{K} w_k \, \mathbf{1}[a_k = a]
\label{eq:hard-voting}
\end{equation}

Donde $\mathbf{1}[\cdot]$ vale $1$ si la condici\'on se cumple y $0$ si
no. Los pesos pueden ser fijos o depender de la confianza de cada
miembro en la decisi\'on actual, $c_k = 1 - H_{\mathrm{norm}}(p_k)$
(Ecuaci\'on~\eqref{eq:shannon-norm}): as\'i, un miembro con una
distribuci\'on m\'as concentrada pesa m\'as.

\subsection{Optimizaci\'on bayesiana de hiperpar\'ametros}

Los agentes de RL tienen hiperpar\'ametros no diferenciables ($\alpha$,
$\gamma$, tama\~no del \emph{buffer}, decaimiento de $\varepsilon$,
anchura de la red, entre otros) que interact\'uan entre s\'i. La
\emph{optimizaci\'on bayesiana} trata el puntaje de validaci\'on
$y = f(x)$ de una configuraci\'on $x$ como una funci\'on de caja negra y
la explora con un modelo probabil\'istico auxiliar. En esta tesis se usa
el estimador estructurado en \'arbol de Parzen (TPE)~\citep{Bergstra2011}
de Optuna~\citep{Akiba2019}. Dado un umbral $y^{*}$, fijado como un
cuantil de los puntajes observados, TPE estima por separado las
densidades $\ell(x) = p(x \mid y < y^{*})$ y
$g(x) = p(x \mid y \geq y^{*})$ con estimadores de densidad por
n\'ucleos, que suavizan las configuraciones observadas. Como el objetivo es maximizar el puntaje, TPE
propone nuevas configuraciones donde el cociente $g(x)/\ell(x)$ es alto,
es decir, donde las configuraciones buenas son m\'as frecuentes que las
malas.

\subsection{Evaluaci\'on de la generalizaci\'on}

Sea $P_{\mathrm{train}}$ la distribuci\'on de las secuencias usadas en
el entrenamiento y $P_{\mathrm{test}}$ la de las secuencias con que se
eval\'ua la pol\'itica. En la condici\'on de referencia ambas
coinciden, $P_{\mathrm{test}} = P_{\mathrm{train}}$. En esta tesis se llama \emph{escenario de generalizaci\'on} a una
condici\'on de evaluaci\'on en la que $P_{\mathrm{test}} \neq
P_{\mathrm{train}}$. En esos escenarios cambia la distribuci\'on de las severidades iniciales, la de las longitudes de secuencia o ambas; la regla de transici\'on (Ecuaci\'on~\eqref{eq:transicion}) y la recompensa no cambian. Los nuevos valores pueden quedar dentro del rango visto en
el entrenamiento (interpolaci\'on) o fuera de \'el
(extrapolaci\'on)~\citep{Packer2018}.

\subsection{Comparaci\'on estad\'istica entre modelos}\label{sec:stats-bg}

Para comparar dos muestras de desempe\~no $A$ y $B$, con $n_A$ y $n_B$
observaciones (en mPES, el desempe\~no de cada secuencia), se usan tres
medidas:

\begin{itemize}
    \item \textbf{Tama\~no del efecto:} la $d$ de Cohen~\citep{Cohen1988}
    expresa la diferencia de medias en unidades de desviaci\'on
    est\'andar:
    \begin{equation}
    d \;=\; \frac{\bar{x}_A - \bar{x}_B}{s_p}
    \qquad
    s_p = \sqrt{\frac{(n_A-1)\,s_A^{2} + (n_B-1)\,s_B^{2}}{n_A+n_B-2}}
    \label{eq:cohen}
    \end{equation}
    Donde $\bar{x}$ y $s^{2}$ son la media y la varianza muestral de cada
    grupo y $s_p$ la desviaci\'on est\'andar combinada. Valores de $|d|$
    cercanos a $0{,}2$, $0{,}5$ y $0{,}8$ se suelen interpretar como
    efectos peque\~no, medio y grande.
    \item \textbf{Significancia:} la prueba $t$ de Welch~\citep{Welch1947}
    contrasta la igualdad de medias sin suponer varianzas iguales:
    \begin{equation}
    t \;=\; \frac{\bar{x}_A - \bar{x}_B}
    {\sqrt{s_A^{2}/n_A + s_B^{2}/n_B}}
    \label{eq:welch}
    \end{equation}
    Los grados de libertad se estiman con la aproximaci\'on de
    Welch--Satterthwaite. El valor $p$ bilateral es la probabilidad de
    observar una diferencia al menos tan grande si las medias fueran
    iguales; un $p$ peque\~no indica que la diferencia no se explica por azar,
    pero no dice cu\'an grande es. En las tablas y figuras se reporta
    $\log_{10} p$.
    \item \textbf{Diferencia entre distribuciones:} la divergencia de
    Kullback--Leibler entre dos distribuciones discretas $p$ y $q$ es:
    \begin{equation}
    D_{\mathrm{KL}}(p \,\|\, q) \;=\;
    \sum_{a} p(a)\log\frac{p(a)}{q(a)}
    \label{eq:kl}
    \end{equation}
    Vale $0$ cuando ambas coinciden y crece con la discrepancia. Se usa
    para comparar la distribuci\'on de acciones de un modelo en un
    escenario con la de la referencia, y, en forma simetrizada
    $\tfrac{1}{2}[D_{\mathrm{KL}}(p\,\|\,q) + D_{\mathrm{KL}}(q\,\|\,p)]$,
    para comparar los histogramas de desempe\~no de dos modelos. Para
    evitar $\log 0$ se suma $\varepsilon = 10^{-9}$ y se renormaliza.
\end{itemize}

\biblio % Needed for referencing to working when compiling individual subfiles - Do not remove
\end{document}
